% GENERATED FILE. Edit canonical lessons, then run npm run generate:books.
% canonical-source-hash: fnv1a64:5fa3c5bbc4444841
% canonical-lessons: TE-C122-jalapatam, TE-C122-divi, TE-C122-avu, TE-C122-gurram, TE-C122-enugu

\chapter{Waterfall, Island, Cow, Horse, Elephant}
\label{ch:te-waterfall-island-cow-horse-elephant}
\hlchaptermodality{122}

\begin{chapteropening}
\textbf{By the end of this chapter:} \emph{I can say a waterfall, an island, a cow, a horse and an elephant.}

Complete the last lesson of chapter 122: i can say a waterfall, an island, a cow, a horse and an elephant.
\end{chapteropening}

\section[jalapātaṁ]{\te{జలపాతం} (jalapātaṁ) — a waterfall}

\label{lesson:TE-C122-jalapatam}



\begin{quote}
Before the new one: say the Telugu for a wave, then the Telugu for a hill.
\end{quote}

\subsection*{You'll want to know: \te{జలపాతం}}
\textbf{\te{జలపాతం}} — \emph{jalapātaṁ} — \textquotedblleft{}a waterfall\textquotedblright{}.

\begin{grammarlens}[title={What you've built}]
One more word for this chapter.
\end{grammarlens}

\subsection*{Your turn}
\begin{itemize}
\raggedright
  \item \emph{Say it:} \emph{jalapātaṁ}
  \item \emph{Say it:} \emph{jalapātaṁ}, once more
  \item \emph{Say it:} say \emph{jalapātaṁ}
  \item \emph{Recall:} say the Telugu for a wave, then the Telugu for a hill, then say \emph{jalapātaṁ} again
\end{itemize}

\begin{tcolorbox}[breakable,colback=teal!4,colframe=teal!35!black,title={Before you move on}]
What does \te{జలపాతం} mean? (\textquotedblleft{}A waterfall\textquotedblright{}.) Say it once more.
\end{tcolorbox}
\section[dīvi]{\te{దీవి} (dīvi) — an island}

\label{lesson:TE-C122-divi}



\begin{quote}
Before the new one: say the Telugu for a hill, then the Telugu for a waterfall.
\end{quote}

\subsection*{You'll want to know: \te{దీవి}}
\textbf{\te{దీవి}} — \emph{dīvi} — \textquotedblleft{}an island\textquotedblright{}.

\begin{grammarlens}[title={What you've built}]
One more word for this chapter.
\end{grammarlens}

\subsection*{Your turn}
\begin{itemize}
\raggedright
  \item \emph{Say it:} \emph{dīvi}
  \item \emph{Say it:} \emph{dīvi}, once more
  \item \emph{Say it:} say \emph{dīvi}
  \item \emph{Recall:} say the Telugu for a hill, then the Telugu for a waterfall, then say \emph{dīvi} again
\end{itemize}

\begin{tcolorbox}[breakable,colback=teal!4,colframe=teal!35!black,title={Before you move on}]
What does \te{దీవి} mean? (\textquotedblleft{}An island\textquotedblright{}.) Say it once more.
\end{tcolorbox}
\section[āvu]{\te{ఆవు} (āvu) — a cow}

\label{lesson:TE-C122-avu}



\begin{quote}
Before the new one: say the Telugu for a waterfall, then the Telugu for an island.
\end{quote}

\subsection*{You'll want to know: \te{ఆవు}}
\textbf{\te{ఆవు}} — \emph{āvu} — \textquotedblleft{}a cow\textquotedblright{}.

\begin{grammarlens}[title={What you've built}]
One more word for this chapter.
\end{grammarlens}

\subsection*{Your turn}
\begin{itemize}
\raggedright
  \item \emph{Say it:} \emph{āvu}
  \item \emph{Say it:} \emph{āvu}, once more
  \item \emph{Say it:} say \emph{āvu}
  \item \emph{Recall:} say the Telugu for a waterfall, then the Telugu for an island, then say \emph{āvu} again
\end{itemize}

\begin{tcolorbox}[breakable,colback=teal!4,colframe=teal!35!black,title={Before you move on}]
What does \te{ఆవు} mean? (\textquotedblleft{}A cow\textquotedblright{}.) Say it once more.
\end{tcolorbox}
\section[gurraṁ]{\te{గుర్రం} (gurraṁ) — a horse}

\label{lesson:TE-C122-gurram}



\begin{quote}
Before the new one: say the Telugu for an island, then the Telugu for a cow.
\end{quote}

\subsection*{You'll want to know: \te{గుర్రం}}
\textbf{\te{గుర్రం}} — \emph{gurraṁ} — \textquotedblleft{}a horse\textquotedblright{}.

\begin{grammarlens}[title={What you've built}]
One more word for this chapter.
\end{grammarlens}

\subsection*{Your turn}
\begin{itemize}
\raggedright
  \item \emph{Say it:} \emph{gurraṁ}
  \item \emph{Say it:} \emph{gurraṁ}, once more
  \item \emph{Say it:} say \emph{gurraṁ}
  \item \emph{Recall:} say the Telugu for an island, then the Telugu for a cow, then say \emph{gurraṁ} again
\end{itemize}

\begin{tcolorbox}[breakable,colback=teal!4,colframe=teal!35!black,title={Before you move on}]
What does \te{గుర్రం} mean? (\textquotedblleft{}A horse\textquotedblright{}.) Say it once more.
\end{tcolorbox}
\section[ēnugu]{\te{ఏనుగు} (ēnugu) — an elephant}

\label{lesson:TE-C122-enugu}



\begin{quote}
Before the new one: say the Telugu for a cow, then the Telugu for a horse.
\end{quote}

\subsection*{You'll want to know: \te{ఏనుగు}}
\textbf{\te{ఏనుగు}} — \emph{ēnugu} — \textquotedblleft{}an elephant\textquotedblright{}.

\begin{grammarlens}[title={What you've built}]
One more word for this chapter.
\end{grammarlens}

\subsection*{Your turn}
\begin{itemize}
\raggedright
  \item \emph{Say it:} \emph{ēnugu}
  \item \emph{Say it:} \emph{ēnugu}, once more
  \item \emph{Say it:} say \emph{ēnugu}
  \item \emph{Recall:} say the Telugu for a cow, then the Telugu for a horse, then say \emph{ēnugu} again
\end{itemize}

\begin{tcolorbox}[breakable,colback=teal!4,colframe=teal!35!black,title={Before you move on}]
What does \te{ఏనుగు} mean? (\textquotedblleft{}An elephant\textquotedblright{}.) Say it once more.
\end{tcolorbox}
